\chapter{Basic Economics}

Why does a concert ticket cost more than a movie ticket? Why do prices go
up a little every year? Why would anyone trade a perfectly good apple for
an orange? Economics answers questions like these. It is a huge subject,
but a few basic ideas explain a surprising amount of how the world works.
We will build them up with a simple story.

\section{Specialization and Trade}

Once upon a time, long before there was money, two people lived in a camp together. They were named Larry and Joan. They both grew oats in nearby fields and mined iron, which they pounded into cookware and other tools for survival.\index{specialization} \index{trade}

Each person needed the tools for both jobs -- plows for oats and shovels and anvils for making iron tools.  And they each had to travel long distances back and forth between their fields and mines.

Finally, one day they decided to \newterm{specialize}: Larry would grow oats and Joan would mine and process iron.  The result was an increase in efficiency: Joan did not need a plow. Larry did not need an anvil. One went to the fields, one went to the mines, and no one needed to travel between. Joan got really skilled at blacksmithing, and Larry got really skilled at raising oats.  With the same number of hours as before specialization, Joan and Larry produced more oats and more iron tools.

As you can imagine,  Larry had lots of oats, but also wanted some iron tools.  Joan had lots of iron tools but also wanted some oats.  So at the end of each week, they would make a trade.  Each person would end up happier because now they each had the oats and the iron tools they wanted.

Here is an important part of economics that many people do not understand: Trade creates wealth.  Before the trades, there was the same amount of oats and iron tools as there would be after the trades, but there was more happiness.  Oats and iron tools are just stuff -- things like happiness, health, and safety are real wealth.

So trade has two purposes:
\begin{itemize}
\item It allows people to specialize, which increases efficiency.
\item It lets people trade what they have for what they want, which increases real wealth.
\end{itemize}

The land and the anvil are both examples of \newterm{capital}: wealth that can be used to make more wealth. Pretty much every productive activity on the planet requires both capital and \newterm{labor}.  Labor is the work of a human. \index{capital} \index{labor}

\section{Money}

As more people arrived in the camp, each person took on a specialization: one grew chickens and traded with eggs, one grew trees and traded with fruit, one mined for gold, etc.

The trading grew more and more complicated: A person who had eggs, but wanted iron tools.  But Joan did not like eggs, she wanted fruit.  So the egg person would try to get fruit but the fruit person wanted oats.  So the egg person would trade eggs for oats, then oats for fruit, and then fruit for iron tools.

Finally gold miner said, "I am going to trade with ANYONE!  I have built a big warehouse, and if you bring me anything useful, I will trade you some gold for it, and store it in my warehouse. And then, if there is anything in my warehouse you want, bring your gold back and I will trade that to you for the gold."

Because everything was stored in the warehouse, they started calling it the \newterm{store}, as in "Take this gold to the store and get me a dozen eggs."

Gold was convenient: It was valuable and easy to carry. It did not spoil. The store made gold coins of different sizes with their weight printed on the front so buying eggs did not involve hacking off and weighing pieces of gold. These gold coins were called \newterm{money} or \newterm{currency}.

\section{Profit}

The store owner tries to buy things with less gold than she sells it for. This difference is called \newterm{profit}. After all, running the store involves a lot of time and some risk -- the fruit might rot before she can sell it.

The most basic idea of profit is this: Each month, a business has to pay some expenses.  Each month that business sells something and generates some revenues.  The revenues minus the expenses is the profit.

Some of the expenses are things like rent and the electricity bill. Some of the expenses are labor costs like the paychecks that the employees got at the end of the month.

\section{Supply and Demand}

If you walked into the store to sell a bushel of oats, the store owner would look in the warehouse: if there was already a lot of oats in there, she would offer you a very small amount of gold for your oats.  If there was not very much, she would offer you more gold for the same bushel of oats.

Of course, the store owner is actually trying to make sure she can sell the oats for at least what she bought them for.  So she is always weighing \newterm{supply} (how many bushels of oats are being produced) with \newterm{demand} (how many bushels of oats people want).

When too many oats are produced, the price drops. As the price drops, more people think "I wanted to eat eggs for breakfast tomorrow, but look how cheap oats are! I think I will eat oats instead."  That is, the demand for oats increases when the price drops.

As the price drops, the farmer who is growing oats thinks "I am doing all this work to grow oats and selling them at that lousy price!? I am going to grow corn".  That is, the supply of oats decreases if the price is very low.

In a free market, supply and demand tend towards equality, and prices are the mechanism by which this occurs. This works astonishingly well. Every day millions of people decide what to buy and what to sell based on the current price.  The result is an almost optimal creation of real wealth with very little waste.

One of the common disasters in economics occurs when the government tries to control prices. For example, people love bread, and sometimes they cannot afford it. A well-intentioned law gets passed that says, ``No one is allowed to sell
bread for more than \$1.00.'' If it costs more than \$1.00 to bake a loaf,
bakers stop baking. Store shelves go empty, and the people the law was
meant to help cannot buy bread at any price.

\section{Competition and Innovation}

Joan, who is the only one in town who makes iron tools, is able to charge a lot for her goods.  After all, she is the only supplier and everyone needs a few iron tools. She has a pile of gold, and it keeps getting bigger.

One of her neighbors notices, and decides to also make iron tools.  He charges less than Joan, and eventually Joan has to lower her prices too. Now when people go to the store for tools, they are much less expensive.

Here is the important point: Competition ensures that businesses cannot charge too much for their products. If you see an industry where profits are unusually high, it indicates that competition is not functioning well and thus the companies in that industry are charging very high prices for their products or services.

The most common form of this is a \newterm{monopoly}.  A monopoly exists when a company has no competition.  In general, a monopoly is very bad for consumers, so most governments have some sort of monopoly prevention laws. In the US, we call these \newterm{antitrust} laws.

Another common way that competition can fail is \newterm{price fixing}. In price fixing, the companies that make a product secretly agree on a
price. For example, say only five companies manufacture an animal feed
additive called lysine.  The leaders of the five companies get together in a hotel room in Paris and decide that instead of competing on price, they will all sell lysine for a high price to ensure they make big profits. This is illegal in most countries. (The lysine example really happened. The companies had to pay \$105 million in fines.)

Notice that competition is inherently inefficient.  For example, Joan had already bought an anvil and dug an iron mine.  For the neighbor to go into the iron tool business, he had to buy an anvil and dig a mine. If it is so inefficient, why does competition seem to create such efficient industries?

The answer is innovation.  Now that Joan and her neighbor are competing, both of them are trying to think of ways to make their products better and cheaper than the other.  They are highly motivated to come up with new ideas and implement them. Competition leads to innovation.

\section{Regulation}

Joan and her neighbor have to purify the iron ore that they pull out of the ground before they make it into tools. The stuff left over is toxic sludge: heavy metals in an acid bath.  Joan and her neighbor pour this out in the yard behind their shops. It seeps into the groundwater, and everyone who drinks from nearby wells gets sick.

Someone tells Joan, "You are making everyone sick; you must be a terrible person! Take your processing byproducts and dump them into a deep hole many miles from town."  Joan replies, "That would take a lot of time and money.  If I do that and my neighbor does not, I will go out of business.  And if I go out of business, my neighbor will double his output and double his waste, and the problem will not be improved at all. I am not a terrible person, but competition prevents me from doing the right thing."

The camp passes a law: ``The byproducts of iron processing must be dumped
into a deep hole at least 10 miles from town.''

Joan and her neighbor both obey the law. They dutifully truck their waste
out to a deep hole far from town.

Iron goods are a little harder to make, so Joan and her neighbor raise their prices some. Everyone who buys iron tools pays the cost of producing them in a safe manner.

The point is this: competition without regulation will lead to a world that none of us would want.

\section{Time Value of Money}

Larry has had a great year: the oats harvest was his biggest ever.  He has a stack of gold coins on his desk when Joan knocks on his door.

"Aluminum, Larry! I want to build a factory to smelt aluminum. Will you lend me four ounces of your gold? I will pay you back in two years, promise. If I do not, you can have my whole iron tool making business."

Larry looks at his gold coins for a long while, and says, "I was planning to build myself a sauna. If I lend you the money, I will need to go two years without a schvitz."

"I will make it worth it.  I will tell you what: Lend me four ounces of gold now and I will pay you back 5!"

Larry took the deal. He used some of his wealth to help buy a productive
asset, hoping to get more wealth back later. This is called
\newterm{investment}.

\begin{Exercise}[title={Larry's interest rate}, label=econ_rate]
Larry lends 4 ounces of gold and gets back 5 ounces two years later. What
annual interest rate, compounded yearly, is that equivalent to?
\end{Exercise}
\begin{Answer}[ref=econ_rate]
We need $4(1 + r)^2 = 5$, so $(1 + r)^2 = 1.25$ and
$1 + r = \sqrt{1.25} \approx 1.118$. That is about 11.8\% per year. (Not
12.5\%: the second year's interest is earned on the first year's
interest too.)
\end{Answer} Nearly all companies involve some capital to get started and to continue growing, so often businesses will seek investment.

\section{Stock}

There are generally two kinds of investment:
\begin{itemize}
\item The return on the investment you make is a predetermined interest rate. For example, if you lent money to a friend to build a factory and said, "Pay me back my principal (the gold I lent you) in five years plus 9\% interest."
\item You own some percentage of the asset.  For example, if your friend was building a factory and could not afford to do it all herself, you might say "I will give you 7 ounces of gold, but then I will own 2\% of your business, and I will take 2\% of all the profits from the business."
\end{itemize}

In the second type of investment, the person making the investment is called a partner or a shareholder.  Companies with a small number of investors are generally partnerships and register with the government as a partnership.

Big companies often need a lot of investors, so they are created as corporations and issue \newterm{shares}.  Each share represents some percentage of ownership.  The percentage is determined by how many shares have been issued by the company.  For example, Google has issued about 12 billion shares. If you buy one share, you own 1 twelve-billionth of Google.

Each quarter, a corporation looks at their bank account and decides whether or not to pay a dividend. For example, a company might decide that every shareholder will get \$1.52 for each share that they own.

\section{Speculation vs. Investment}

When you invest your money in a productive asset (like a factory), your asset creates jobs and wealth.

When you buy an unproductive asset (like bitcoin), your asset creates no jobs or wealth.  You are just hoping that eventually you can sell it for more than you paid.

Every year, economists try to estimate the \newterm{gross domestic product} (or GDP) of every nation on the planet.  GDP is the value of all products and services produced in a country in a year. Naturally, large countries tend to have larger GDPs than small nations, so we often compare per-capita GDP -- that is, the GDP of a nation divided by the number of people in the nation.  It is not a perfect measure, but generally nations with a large per-capita GDP feel wealthier than nations with a small per-capita GDP.

Here are the approximate per-capita GDPs of a few nations around 2023:
\begin{itemize}
\item Monaco (the highest of any nation): \$240,000
\item Ireland : \$103,000
\item United States: \$81,000
\item South Korea: \$33,000
\item Hungary: \$22,000
\item Russia: \$13,000
\item China: \$13,000
\item Afghanistan: \$400
\item Burundi (the lowest of any nation): \$200
\end{itemize}

Does speculation do any social good? You can argue that speculators, who are always trying to find a mispriced asset, are correcting flaws in the price: If an asset has a price that seems to be too low, the speculator will buy some as soon as possible, and this will raise the price to where it should be. Thus speculation can be part of the magic that maintains the balance between price, supply, and demand.

In general, governments are much more interested in investment than speculation.

\section{Taxes}

The government provides services to its citizens.  These services cost money, so the government demands that its citizens pay taxes.

Some governments provide more services to their citizens. These are called "Big Government" nations and tend to have taxes of more than 30\% of GDP. Germany's taxes are nearly 40\% of its GDP.

Some governments provide fewer services to their citizens. These are called "Small Government" nations and tend to have taxes of less than 25\% of GDP. Mexico's taxes are about 17\% of its GDP.

The US is somewhere in the middle with taxes being about 28\% of GDP.

Most governments try to make their tax code "progressive" not "regressive."  Each person pays a percentage of their income in taxes. With a progressive tax system poor people pay a smaller percentage of their income -- giving them more to survive on.  Rich people pay a higher percentage of their income to make up for the tax breaks that were given to the poor people.

There are several ways of doing taxation, and most nations use some sort of combination to make it difficult for people to avoid paying their fair share:
\begin{itemize}
\item Income tax: You made \$$X$ dollars this year? Send us $y$ percent of that.
\item Property tax: Your house is worth \$$X$? Send us $y$ percent of that.
\item Sales tax: You bought something for \$$X$? Send us $y$ percent of that.
\item Corporate tax: The company made a profit of \$$X$? Send us $y$ percent of that.
\item Inheritance tax: Your dad died with \$$X$ in the bank? Send us $y$ percent of your inheritance.
\item Capital gains tax: Your investment made you a profit \$$X$? Send us $y$ percent of that profit.
\item Wealth tax: You have \$$X$ worth of assets? Send us $y$ percent of that.
\item Tariffs: You, a foreign company, want to sell a product in the US? Send us \$$y$ for each one that you export to the US.
\end{itemize}

In the US, counties and cities often collect property taxes and sales taxes.  States often collect sales taxes and income taxes.  The federal government collects income taxes, corporate taxes, inheritance taxes, and capital gains taxes. At this time, we do not have any wealth taxes in the US.

No one likes taxes, but we rely upon the government for our safety, our infrastructure, and for some sort of safety net if we are unable to work for some reason.

The tax code often includes incentives for behavior that the government would like to promote. For example, some governments might give you a tax break if you spend money to improve the insulation on your house.

Sales taxes are generally considered regressive because poor people tend to spend a much higher percentage of their income on consumer goods. Rich people save and invest a considerable percentage of their income.

\section{Tariffs}

Tariffs sound like a great idea: Let us get people in other nations to pay our taxes for us! And if we make the tariffs high enough, companies in other countries will not be able to compete with ours! In practice, tariffs seldom work out well.

Let us work through an example using simple numbers. Here are the players:
\begin{itemize}
\item Walmart sells T-shirts in the US. They will buy the cheapest shirt they can and sell them for 50\% more than the price they purchased them for.
\item Rex Shirts is an American company that can produce T-shirts for \$18. For the business to be worth doing, they need to make a \$2.00 profit on each shirt. So they can sell the T-shirts for \$20.
\item Dhaka Shirts is a T-shirt manufacturer in Bangladesh.  It can produce a T-shirt for \$9. For the business to be worth doing, they need to make a \$1.00 profit on each shirt. So they can sell the T-shirts for \$10.
\item Jenna wants to buy a T-shirt for her son. She lives near a Walmart.
\end{itemize}

\subsection{No Tariffs}

If there are no tariffs on T-shirts from Bangladesh, Walmart buys one for \$10 from Dhaka Shirts and sells it to Jenna for \$15.

Note the two great things that happened here:
\begin{itemize}
\item Jenna wanted the T-shirt more than she wanted \$15, so wealth was created simply by the fact that trade occurred.
\item A company in Bangladesh got \$10. Some of that went to poor people working on sewing machines, some went to the cotton suppliers, and some was profit for the company's owners. Those people paid taxes to the Bangladeshi government. When they think about invading their neighbor Myanmar, they know that the war will end the business relationship and everyone will be poorer as a result. Peace is more likely.
\end{itemize}

Note two bad things that happened here:
\begin{itemize}
\item The US government got no tax revenues when the shirt came into our country.
\item Rex Shirts was unable to compete with Dhaka Shirts on this product. They will need to find a product where they can compete, or innovate, or go out of business.
\end{itemize}

\subsection{Low Tariff}

Let us say that the US puts a \$5 tariff on T-shirts from Bangladesh: Walmart buys one for \$15, \$10 goes to the Dhaka Shirts (still cheaper than Rex Shirts) and \$5 goes to the US government. Walmart tries to sell the shirt to Jenna for \$22.50.

Notice that the price to Jenna went from \$15 to \$22.50 -- a 50\% increase.  If she buys it, the trade still occurs and the government gets \$5.  This is why tariffs are sometimes compared to sales taxes: Jenna pays more and most of the difference goes to the US government.  And, like sales taxes, it is regressive. Poor people spend a larger percentage of their income on T-shirts than rich people.

\subsection{High Tariff}

When the tariffs are being created, Rex Shirts might send their lobbyist to Washington.  Maybe he can convince the government to put a \$15 tariff on the shirt. Now the Dhaka Shirts Tee costs Walmart \$25. So Walmart calls up Rex Shirts and asks for a quote.

Rex Shirts could sell the shirt for \$20, but they know that the competition cannot go lower than \$25.  So Rex Shirts says, "We can sell you a T-shirt for \$24."  Walmart adds their margin, and tries to sell the shirt for \$36.

Now the price has gone from \$15 to \$36. If Jenna buys it, the employees
at Rex Shirts get a chance to make a shirt, and Rex himself makes a
spectacular profit of \$6. But no shirt crossed the border, so the
government collects nothing at all. And the price for Jenna has increased
by 140\%. The tariff's real effect was to move money from Jenna's pocket to
Rex's.

\subsection{Lost Trades}

Here is a more subtle point: any tariff makes it less likely that Jenna
will buy the shirt at all. And every trade creates a little wealth: Jenna
wants a T-shirt more than she wants \$15, and Dhaka Shirts wants \$10 more
than it wants the shirt. If the tariff stops this trade, that wealth is
never created. Economists call wealth lost this way \newterm{deadweight
loss}\index{deadweight loss}.

In a healthy economy, money changes hands constantly, and each trade
creates a little wealth. Every barrier to trading means some of those
trades never happen.

\section{Libertarianism vs Socialism}

You can imagine that there is an axis: at one end is true libertarianism: the free market determines everything; there are no regulations.  At the other end is true socialism: the government owns every productive asset and sets all prices; people are just consumers and labor.

Neither extreme works:
\begin{itemize}
\item Libertarian: As discussed before, without regulation, competition drives companies to be more and more irresponsible. Everyone creating their own infrastructure is wildly inefficient.
\item Socialist: Setting prices and allocating capital into investments is very difficult; it is best done by someone who will actually lose their savings if they invest it badly. Also, when all the money and power flows through the government, it is difficult for government employees to resist corruption.
\end{itemize}

Most countries are somewhere in the middle: France is a little more socialist than the US. China is more socialist than France. Somalia is more libertarian than the US.

\section{The Money Supply}

In the earlier example, we used gold coins as the money. Gold coins are heavy, and it is easy to dope them with less valuable metals. Eventually, governments around the world issued paper bank notes that
stood for gold but were easier to work with. If you really wanted to, you could take your bank note to the government and it would give you the appropriate amount of gold. For about a century, a US dollar could be traded in for about
$\frac{1}{20}$ of an ounce of gold.

Sometimes the value of gold goes up.  Maybe at the beginning of the year you could buy a bushel of oats for 1/20 of an ounce of gold, and at the end of the year you could buy two bushels of oats for the same amount of gold. So holding dollars was equivalent to speculating on gold, and a lot of people did this.

Remember that the government does not really want you to speculate on
gold; it is not a productive asset. It does not want you to stash dollars
under your bed, either. It wants you to invest them in productive assets.
The economy is healthier if the value of the dollar decreases a little
every year. That gives you a reason to get your dollars out from under the
bed and invest them in something productive.

How do you make dollars decrease in value? First, you unhitch them from gold. Then, anytime it looks like the value of the dollar might increase, you print a few more and get them into circulation.

In the US, the Federal Reserve System was created in 1913 to control the money supply, but the ability to convert dollars to gold was not fully terminated until 1971.

The Federal Reserve (or just ``the Fed'') steers the economy mainly by
controlling the interest rate that banks pay to borrow money. When it wants
more money in the system, it lowers that rate, and loans get cheaper. This tends to juice the economy a little as businesses borrow money to create productive assets. However, it also causes the value of each dollar to decrease slightly; we call this \newterm{inflation}.  A little inflation, say 2 or 3 percent, is good for the economy. If it is more, it is difficult for wage growth to keep up. People who live paycheck-to-paycheck are hurt the most by high inflation.

Of course, every president wants the Fed to offer money at low interest rates during their term because this gives a short-term boost in growth. This is why the Fed is designed to be independent: the president and
Congress cannot order it to change interest rates.  The Fed is led by very boring economists who try to figure out how fast to move money into the economy so that we get nice, consistent growth over the decades.

\section{Labor vs Capital}

Most people make most of their money by being labor -- they get a paycheck from a corporation.  Some people make money from the businesses that they have invested in -- they make money by owning things. For our purposes, we will call them "owners."

Owners tend to have a lot of control over things. They are the rich and have the ear of the government. Over the years, the government has passed many laws that primarily help owners. Also, technology tends to make labor even less powerful. In the US, we have seen an increasing gap between the income and wealth of the owners and income and wealth of labor. This should be concerning to all of us, and we should think about it whenever we consider:
\begin{itemize}
\item Changes to the tax code
\item Raising the minimum wage
\item Laws concerning labor unions
\item Worker safety regulations
\end{itemize}

Here is the other concerning thing: It has become harder to change groups.  If your parents are owners, you will probably also be an owner all your life.  If your parents are labor, you will probably be labor all your life.  There are exceptions, of course, but changing groups was a lot more common 50 years ago than it is today. Here are two important parts to fixing this:
\begin{itemize}
\item Quality public schools -- opportunity starts with a top-notch education. Every child should get a free high-quality education until they are an adult. Public universities should be affordable and excellent. When someone says "School Choice," they are really saying "I am not interested in creating opportunities for other people's children."
\item The Inheritance Tax -- The inheritance tax prevents dynasties of owners, and ensures that wealthy families help pay for the government services that everyone relies upon.
\end{itemize}

\section{Capitalism Works}

Among young idealists there is a dream of eliminating the owning class altogether. In their dream, everyone is both an owner and labor. It is an appealing idea, and one that will never work.

One of the most difficult problems is capital allocation -- where do we invest money today such that it will generate wealth in the future? When you see the demolition of huge new apartment buildings in China, you are seeing the result of bad capital allocation.  Capital allocation is best done by people who will win or lose depending on how good the investment is -- capital allocation by committee is a recipe for economic disaster.

You need every job filled: we need roofers to go up on the roof in the summer heat. Who will do that if we have made failing impossible?

You need to make it possible to win big: Innovation needs to deliver more than a paycheck. Taking a gamble on creating a new business needs to result in significant wealth if it goes well.

Every nation that has tried to eliminate winners and losers has ended up with a stagnant economy. It just does not work.

\begin{Exercise}[title={Tariff and Price}, label=econ1]
In the tariff example, Walmart pays \$10 for a shirt and sells it for
50\% more. Suppose the tariff were \$8 instead of \$5. What would Jenna pay,
and by what percent would her price rise compared with no tariff? How much
would the government collect per shirt?
\end{Exercise}
\begin{Answer}[ref=econ1]
Walmart's cost is $\$10 + \$8 = \$18$, still less than Rex Shirts' \$20,
so Walmart keeps buying from Dhaka Shirts. Jenna pays
$1.5 \times \$18 = \$27$, up from \$15, an increase of
$\frac{12}{15} = 80\%$. The government collects \$8 per shirt.
\end{Answer}
